\section{Recognizing the Grassmannian}\label{sec:grassmannian}

We finish the pseudoindex-five case by interpreting the numerical constraint
on the fibres of the evaluation map.  Throughout this section, $X$ is a
smooth complex Fano sixfold with nef tangent bundle, Picard number one and
pseudoindex five.  We use its minimal rational component $V$, universal
family $\pi\colon U\to V$, and evaluation $e\colon U\to X$.
Thus $e$ is smooth with connected three-dimensional fibres. For a general
point $x\in X$, let $F=e^{-1}(x)$ and let $L$ be the tangent-map
polarization of Lemma~\ref{lem:marked-tangent}.
Our goal is to recover $X$ from the projective geometry of the tangent
directions to its minimal rational curves.
For this polarization, Lemma~\ref{lem:finite-ratio} gives
\begin{equation}\label{eq:vmrt-ratio-input}
 32-20u+3u^2=0,
 \qquad u=\frac{c_1(T_F)\cdot L^2}{L^3}.
\end{equation}
We first show that the tangent image spans $\mathbb P^5$. Its degree
bound and the sectional-genus formula will then determine $L^3$ and
the degree of the tangent map. Finally, normality of the image will
let us use the smoothness of $F$ to select the Segre threefold among
the varieties of minimal degree.

\begin{lemma}\label{lem:vmrt-nondegenerate}
For a general point $x\in X$, the threefold
$\mathcal C_x\subset\mathbb P(T_x^*X)\simeq\mathbb P^5$ is linearly
nondegenerate.
\end{lemma}

\begin{proof}
Write $\widehat{\mathcal C}_x\subset T_xX$ for its affine cone and
$W_x=\operatorname{Span}(\widehat{\mathcal C}_x)$.  Suppose that $W_x$
is proper.  Since $\mathcal C_x$ has dimension three,
$4\leq\dim W_x\leq5$.  On a dense open subset of $X$, these spans
define the distribution $W\subset T_X$.

We use two results of Hwang and Mok
\cite[Propositions~1.2.1 and~1.2.2]{HwangMok1999}.  For an irreducible
variety of minimal rational tangents on a uniruled projective manifold
with $b_2=1$, a proper spanned distribution cannot be integrable.  On
the other hand, it is integrable if the bivectors
\begin{equation}\label{eq:tangent-bivectors}
 \alpha\wedge v,\qquad
 \alpha\in\widehat{\mathcal C}_x\setminus\{0\}
 \text{ a smooth point},\quad
 v\in T_\alpha\widehat{\mathcal C}_x,
\end{equation}
span $\bigwedge^2W_x$ at general $x$.  These results apply here:
$b_2(X)=1$ was established earlier, and $\mathcal C_x$ is irreducible
as the image of $F$.

We verify the spanning condition directly, without assuming that
$\mathcal C_x$ is smooth.  Otherwise a nonzero alternating form
$\omega\in\bigwedge^2W_x^*$ annihilates all bivectors
\eqref{eq:tangent-bivectors}.  The one-form $\omega(\alpha,d\alpha)$
vanishes on the smooth cone locus.  Its exterior derivative is the
restriction of $2\omega$, so every four-dimensional tangent space
$T_\alpha\widehat{\mathcal C}_x$ is isotropic.  An alternating form
of rank $2r$ on an $m$-dimensional space has isotropic subspaces of
dimension at most $m-r$.  Since $m\leq5$ and $r\geq1$, we must have
\[
 4\leq m-r\leq4,
\]
and hence $m=5$ and $r=1$.

Projection to the two-dimensional symplectic quotient
$W_x/\ker\omega$ has differential of rank at most one on the smooth
cone locus.  In characteristic zero, the Zariski closure of the image therefore
has dimension at most one.  It is an irreducible cone, hence is either
the origin or a line.  But $\widehat{\mathcal C}_x$ spans $W_x$, so
its projection spans the two-dimensional quotient, a contradiction.
Thus the bivectors do span $\bigwedge^2W_x$, making $W$ integrable
and contradicting the other Hwang--Mok result.
\end{proof}

\begin{proposition}\label{prop:pseudoindex-five-recognition}
Every smooth complex Fano sixfold with nef tangent bundle, Picard number
one and pseudoindex five is isomorphic to $\operatorname{Gr}(2,5)$.
\end{proposition}

\begin{proof}
The geometric constraints~\eqref{eq:N-products} and
Lemma~\ref{lem:finite-ratio} give $u=4$ or $u=8/3$. We first determine
the degree of $L$.  By Bertini, two general members of $|L|$ meet in a smooth
curve $B$.  Both successive divisors are connected:
for a smooth projective variety $Z$ of dimension at least two and an
ample divisor $D$, Kodaira vanishing and Serre duality give
$H^1(Z,\mathcal O_Z(-D))=0$, and the divisor exact sequence then gives
$H^0(D,\mathcal O_D)=\mathbb C$.  Applying this first to $F$ and then
to its smooth ample surface section proves that $B$ is connected.
Adjunction gives
\begin{equation}\label{eq:vmrt-sectional-genus}
 g(B)=1+\frac{(K_F+2L)\cdot L^2}{2}
     =1+\left(1-\frac u2\right)L^3.
\end{equation}
Here and below $L^3$ denotes its top intersection number.

By Lemma~\ref{lem:vmrt-nondegenerate}, $\mathcal C_x$ is a
nondegenerate threefold in $\mathbb P^5$, so its degree is at least
three.  The degree formula for the finite morphism $\tau_x$ yields
\[
 L^3=\deg(\tau_x)\deg(\mathcal C_x)\geq3.
\]
If $u=4$, Equation~\eqref{eq:vmrt-sectional-genus} gives
$g(B)=1-L^3<0$.  Thus $u=8/3$, and the same equation gives
$g(B)=1-L^3/3\geq0$.  It follows that
\[
 L^3=3,\qquad\deg(\tau_x)=1,\qquad\deg(\mathcal C_x)=3.
\]

The del Pezzo--Bertini classification of varieties of minimal degree
\cite[Theorem~0.1]{EisenbudGreenHulekPopescu2006} now makes
$\mathcal C_x$ a rational normal scroll, possibly a cone.  Indeed,
linear spaces and quadrics have smaller degree, while a cone over the
Veronese surface has degree four.  A three-dimensional scroll of
degree three has parameters $a_1,a_2,a_3\geq0$ with
$a_1+a_2+a_3=3$.

To turn the finite degree-one tangent map into an isomorphism, we need
normality of its target. These scrolls are normal, including the cone
cases, as follows. Their homogeneous coordinate rings are generated
by the monomials
\[
 z_i s_0^j s_1^{a_i-j},\qquad
 1\leq i\leq3,\quad0\leq j\leq a_i.
\]
These monomials generate exactly the invariant subring of
$\mathbb C[z_1,z_2,z_3,s_0,s_1]$ for the multiplicative-group action
with weights $(-a_1,-a_2,-a_3,1,1)$, graded by total degree in the
$z_i$. In this grading every displayed generator has degree one.
To see that they generate the invariant ring, distribute the
$s_0$ and $s_1$ exponents of an invariant monomial among its $z_i$
factors, each of which requires exactly $a_i$ such exponents.  An
invariant subring of a normal domain is normal: an element of its
fraction field integral over it belongs to the original domain and
remains invariant.  Thus the scroll is normal.

The finite birational morphism $\tau_x\colon F\to\mathcal C_x$ is
therefore an isomorphism.  In particular, $\mathcal C_x$ is smooth.
If some $a_i$ vanished, the scroll would be a nontrivial cone over a
nonlinear variety and hence singular along its vertex.  Consequently
all $a_i$ are positive, so $(a_1,a_2,a_3)=(1,1,1)$.  We have proved
the projective equivalence
\[
 \mathcal C_x\simeq
 \mathbb P^1\times\mathbb P^2\hookrightarrow\mathbb P^5,
 \qquad\text{embedded by }\mathcal O(1,1).
\]

This is the variety of minimal rational tangents of
$S=\operatorname{Gr}(2,5)$.  At a two-plane $A\subset\mathbb C^5$,
$T_AS=\operatorname{Hom}(A,\mathbb C^5/A)$, and tangent directions
to lines are exactly the rank-one maps.  Their projectivization is
the displayed Segre variety.  For example, in graph coordinates the
Pl\"ucker coordinates consist of the constant coordinate $1$, the entries of a map, and its
$2\times2$ minors; a direction determines a line precisely when the
quadratic minors vanish.

Finally, Mok's recognition theorem
\cite[Main Theorem, \S2.1]{Mok2008Recognition} states that a Fano
manifold of Picard number one whose general variety of minimal
rational tangents is projectively equivalent to that of an irreducible
compact Hermitian symmetric space is isomorphic to that space.
The model $S=\operatorname{Gr}(2,5)$ has rank two, and we have
verified the required projective equivalence for the minimal component
$V$.  The theorem gives $X\simeq S$.
\end{proof}
